
\chapter{버전별 변경사항}

\newversionsection{0.6.6}{2026-10-06}{
    \par -
}{
    \par -
}{
    \begin{itemize}
        \item 형상 생성 단계에서의 랜덤함수 개입으로 시뮬레이션 시마다 결과값이 달라지던 문제 수정.
        \item 동일 실 내에 면적비율이 0.1\% 미만인 바닥면이 존재할 경우 바닥난방시스템 적용시 유량 분배 과정에서 발생하던 오류 수정
        \item 냉동기 모델 생성 후 ExpandObjects 단계에서 발생하던 오류 수정
    \end{itemize}
}{
    \par -
}

\newversionsection{0.6.5}{2026-08-07}{
    \par -
}{
    \par -
}{
    \begin{itemize}
        \item 히트펌프의 \texttt{fuel\_type} 입력이
              EnergyPlus의 VRF 생산설비 객체에 반영되지 않던 문제 수정.
        \item 흡수식 냉동기의 \texttt{fuel\_type} 입력이
              연동 보일러에 반영되지 않고 천연가스로 고정되던 문제 수정.
        \item 흡수식 냉동기의 \texttt{cop\_cooling} 입력이
              연료사용량 계산용 성능계수에 반영되지 않던 문제 수정.
        \item 베네시안 블라인드를 \texttt{InteriorShade}로 잘못 생성하던
              문제를 수정하고, \texttt{WindowMaterial:Blind}의
              슬랫 간격 필드명 및 가시광선 투과율 필드를 수정.
        \item Optional 속성인 창의 \texttt{blind}와 방열기의
              \texttt{capacity\_heating}을 생략하면
              GRM 파싱 중 오류가 발생하던 문제 수정.
        \item 냉동기 압축기 종류를 문자열로 전달한 경우
              \texttt{to\_idf\_curve\_object} 호출 과정에서
              오류가 발생하던 문제 수정.
    \end{itemize}
}{
    \par -
}

\newversionsection{0.6.4}{2026-07-27}{
    \par -
}{
    \begin{itemize}
        \item 지면(ground), 단열면(adiabatic)을 경계조건으로 가지는 면(surface)에 개구부를 입력하였을 때 명시적 에러 발생
    \end{itemize}
}{
    \begin{itemize}
        \item 문자열 입력값에 ,;! 가 포함되었을 때 오류 발생하던 문제 수정
        \item 입력값의 범위 검증 시 경계값을 포함하지 않는 경우에도 경계값 포함으로 검증되던 오류를 수정
    \end{itemize}
}{
    \par -
}

\newversionsection{0.6.3}{2026-07-19}{
    \begin{itemize}
        \item 내장 런처가 grm 파일 실행을 지원
        \item 내장 런처가 진행 상황을 표시 (디버그, 시뮬레이션 중, 결과 처리 중 3단계별 완성 여부를 표시)
    \end{itemize}
}{
    \begin{itemize}
        \item 용도프로필을 EnergyPlus 스케줄로 변환하는 알고리즘의 전면 수정 (Technical Reference Manucal 참고, 입력 가능한 용도프로필의 종류 및 명칭은 기존과 동일함)
        \item 자연환기 및 기계환기(전열교환기 적용 시)에서 인당 환기 유량을 2L/s에서 8.3L/s로 변경
        \item 태양광 패널의 발전량 산정 시 현장에서 사용되지 않고 남는 전력은 제외하도록 변경
    \end{itemize}
}{
    \begin{itemize}
        \item 급탕이 아닌 지역난방 연료 사용 시 결과에 집계되지 않던 오류 수정
        \item 허가일자가 법규 시행일과 같을 때 이전 법규가 적용되던 오류 수정
        \item 베네시안(venetian) 블라인드 적용 시 오류 발생하던 문제 수정
        \item 전기방열기의 용량(선택 옵션)을 미입력 시 실행이 불가하던 오류 수정
    \end{itemize}
}{
    \par -
}

\newversionsection{0.6.2}{2026-05-27}{
    \par -
}{
    \begin{itemize}
        \item \inlinecode{EPlusSimpleCLI.exe}가 기존 \inlinecode{runEngine.bat}의 CLI 명령어 환경을 대체함.
        \item \inlinecode{EPlusSimpleLauncher.exe}가 기존 \inlinecode{runLauncher.bat}의 런처 실행을 대체함.
        \item 런처는 별도의 console과 Web GUI를 실행하지 않고, 별도 창으로 실행됨.
    \end{itemize}
}{
    \begin{itemize}
        \item 벽체 면적이 또는 바닥면적이 아주 작은 경우 실내온도가 발산하던 오류 수정.
    \end{itemize}
}{
    \begin{itemize}
        \item 더이상 batch 스크립트를 활용하여 CLI 명령어로 엔진 실행 또는 Launcher 실행할 수 없음.
    \end{itemize}
}

\newversionsection{0.6.1}{2026-01-27}{
    \begin{itemize}
        \item 침기를 optional 변수로 변경 (\textit{null}로 입력 시 외기와 맞닿은 창문이 있는 경우 $1.5ACH50$, 없는 경우 $0ACH50$로 자동 입력).
        \item ``교실(어린이집)'' 용도프로필 추가 (방학이 없는 ``교실(초중고)'' 프로필과 동일함)
    \end{itemize}
}{
    \begin{itemize}
        \item 건물 내벽의 구조체를 \textit{null}로 입력 시 외벽의 법규기준 열관류율에 준해 자동 입력되던 것에서, 석고보드($12.5mm$)--글라스울($50mm$)--석고보드($12.5mm$)로 고정 입력되도록 변경.
    \end{itemize}
}{
    \begin{itemize}
        \item 태양광 발전 시스템에서 패널 면적이 $1/4$배로 작게 계산되던 오류 수정.
        \item grr파일의 ``summary\_per\_area'' 및 ``summary\_gross'' 항목에서 연료별 합계에 태양광 발전량도 포함되는 오류 수정 (발전량은 차감하여 출력하도록 수정).
    \end{itemize}
}{
    \par -
}


\newversionsection{0.6.0}{2025-12-08}{
    \begin{itemize}
        \item Regression Test Report 문서 추가.
        \item Release Note 문서 추가.
    \end{itemize}
}{
    \par-
}{
    \par-
}{
    \begin{itemize}
        \item 기본으로 제공되는 용도프로필 목록에서 ``구내식당(초중고)'', ``주방 및 조리실 (초중고)'', ``체육시설 (초중고)'' 삭제.
        \item ``custom'' 프로필 입력은 이제 허용되지 않음 (\texttt{zone}의 속성 중 ``profile\_id'' 삭제, grm 구조에서 \texttt{profile} 및 \texttt{profile\_components} 아래의 항목들 모두 삭제).
    \end{itemize}
}


\newversionsection{0.5.2}{2025-11-28}{
    \begin{itemize}
        \item 수정된 ASHRAE 140 케이스 600 모델을 예시파일로 추가.
        \item grr 결과파일 중 \textit{summary\_per\_area}, \textit{summary\_gross} 항목에 연료별, 용도별 합계를 추가
    \end{itemize}
}{
    \par -
}{
    \par -
}{
    \par -
}

\newversionsection{0.5.1}{2025-11-17}{
    \begin{itemize}
        \item 공급설비의 종류로 전기바닥난방(electric\_radiant\_floor)을 추가 (electric\_radiator와 동일하게, 이름 외 별도 속성 변수 없음).
    \end{itemize}
}{
    \begin{itemize}
        \item 냉난방 설비의 자동 용량 산정 (Autosizing) 기간을 1월 15일부터 8월 15일까지에서 1월 한달간 및 8월 한달간으로 변경. 이로 인해 용량 미입력시 자동 계산되던 설비의 용량이 변화할 수 있음.
        \item 패키지에어컨이 EnergyPlus의 AirConditioner:VariableRefrigerantFlow 객체를 사용 (기존에는 ZoneHVAC:WindowAirConditioner사용; 공조기+EHP와 동일한 객체 사용하도록 변경).
        \item 패키지에어컨 및 EHP 사용시 냉방 코일의 현열비 및 정격 풍량을 자동 산정하지 않고 각각 0.7, 10L/$s\cdot m^2$으로 고정.
    \end{itemize}
}{
    \begin{itemize}
        \item 패키지에어컨 사용시 냉방에너지가 계산되지 않던 (EnergyPlus 시뮬레이션에서 항상 가동하지 않던) 오류 수정.
        \item EHP 사용시 특정 조건에서 발생하던 EnergyPlus 실행 오류 수정.
        \item 난방유 사용시 grr 파일에서 결과에 집계되지 않던 오류 수정.
        \item EnergyPlus 실행 후 err 파일을 읽을 수 없던 경우 발생하던 오류를 EnergyPlusError Exception으로 예외처리 가능하도록 수정.
    \end{itemize}
}{
    \par - 
}
*\textit{V0.5.1의 수정사항은 V0.5.0을 기준으로 작성되었습니다.}
